\documentclass[12pt, twoside]{article}
%\documentclass[12pt, twoside]{article}
\usepackage[letterpaper, margin=1in, headsep=0.2in]{geometry}
\setlength{\headheight}{0.6in}
%\usepackage[english]{babel}
\usepackage[utf8]{inputenc}
\usepackage{microtype}
\usepackage{amsmath}
\usepackage{amssymb}
%\usepackage{amsfonts}
\usepackage[nomessages]{fp} %\FPeval{\var-name}{2*sin(pi/6)}
\usepackage{siunitx} %units in math. eg 20\milli\meter
\usepackage{yhmath} % for arcs, overparenth command
\usepackage{tikz} %graphics
\usetikzlibrary{quotes, angles, arrows, arrows.meta}
\usepackage{pgfplots}
\usepgfplotslibrary{fillbetween}
\pgfplotsset{compat=1.16}
\usepackage{graphicx} %consider setting \graphicspath{{images/}}
\usepackage{parskip} %no paragraph indent
\usepackage{enumitem}
\usepackage{multicol}
\usepackage{venndiagram}

\usepackage{fancyhdr}
\pagestyle{fancy}
\fancyhf{}
\renewcommand{\headrulewidth}{0pt} % disable the underline of the header
\raggedbottom
\hfuzz=2mm %suppresses overfull box warnings

\usepackage{hyperref}

\fancyhead[LE]{\thepage}
\fancyhead[RO]{Markscheme}
\fancyhead[LO]{La Scuola d'Italia / Dr.\ Huson / IB Math 11 \\* 30 September 2026}

\begin{document}
\subsubsection*{1.7 Classwork: Financial calculations and compound interest --- Markscheme}

\begin{enumerate}[itemsep=0.3cm]

%--- Problem 1: Simple interest as an arithmetic sequence [6 marks] ---
\item[1.] [Maximum mark: 6]

Maria invests \$$2000$ at $3.5\%$ per year simple interest.

\medskip

\textbf{(a)} $I = 2000 \times \dfrac{3.5}{100} = $ \$$70$ \hfill \textbf{A1}

\emph{A value of \$$7000$ (using $3.5$ in place of $0.035$) earns no
mark.}

\medskip

\textbf{(b)} \$$2070$, \ \$$2140$, \ \$$2210$ \hfill \textbf{A1}

\emph{All three values required for the mark.}

\emph{FT: award A1 for $2000 + I$, $2000 + 2I$, $2000 + 3I$ with the
candidate's own $I$ from (a).}

\medskip

\textbf{(c)} $V_{n} = 2000 + 70n$ \hfill \textbf{A1}

\emph{Accept $V_{n} = 2070 + 70(n - 1)$ or any equivalent form.}

\emph{$V_{n} = 2000 + 70(n - 1)$ earns no mark: it gives \$$2000$ at
the end of year $1$.}

\emph{FT: award A1 for $2000 + In$ with the candidate's own $I$ from
(a).}

\medskip

\textbf{(d)} $V_{12} = 2000 + 70 \times 12 = $ \$$2840$ \hfill \textbf{A1}

\emph{FT: award A1 for the value obtained by substituting $n = 12$
into the candidate's own expression from (c).}

\medskip

\textbf{(e)} $2000 + 70n > 3000$ \hfill \textbf{M1}

$V_{14} = $ \$$2980$ \ ($< 3000$) \quad and \quad $V_{15} = $ \$$3050$
\ ($> 3000$)

$15$ years \hfill \textbf{A1}

\emph{Accept solving the inequality: $n > 14.28\ldots$, so $15$
years.}

\emph{Accept a table search: full marks for showing both adjacent
years and stating $15$. A single row (after $15$ years the value is
\$$3050$, so $15$) still earns full marks, with the feedback note that
the adjacent year $14$ is what justifies ``first''. See
markscheme-spec §18.}

\emph{$14.3$ is not an acceptable final answer --- the answer must be
a whole number of years. Rounding down to $14$ earns the M1 only.}

\emph{FT: award A1 for the smallest whole number of years correctly
obtained from the candidate's own expression in (c).}

\newpage

%--- Problem 2: Compound interest, annual vs monthly, doubling time [7 marks] ---
\item[2.] [Maximum mark: 7]

\$$2000$ invested. Bank A: $3.5\%$ compounded annually. Bank B:
$3.4\%$ compounded monthly.

\medskip

\textbf{(a)} $2000 \times 1.035^{12}$ \hfill \textbf{M1}

$=$ \$$3022.137\ldots \approx$ \$$3022.14$ \hfill \textbf{A1}

\emph{Accept \$$3020$ (3 s.f.), \$$3022$ or \$$3022.14$. Accept the
TVM solver: $N = 12$, $I\% = 3.5$, $PV = -2000$, $PMT = 0$,
$P/Y = C/Y = 1$.}

\emph{A multiplier of $0.035$ or $1.35$ earns no marks. An exponent of
$11$ (giving \$$2919.94$) earns the M1 only.}

\emph{Simple-interest working ($2000 + 12 \times 70 = 2840$) earns no
marks.}

\medskip

\textbf{(b)} $2000\left(1 + \dfrac{3.4}{100 \times 12}\right)^{12
\times 12}$ \hfill \textbf{M1}

$=$ \$$3005.879\ldots \approx$ \$$3005.88$ \hfill \textbf{A1}

\emph{Accept \$$3010$ (3 s.f.), \$$3006$ or \$$3005.88$. Accept the
TVM solver: $N = 144$, $I\% = 3.4$, $PV = -2000$, $PMT = 0$,
$P/Y = C/Y = 12$.}

\emph{M1 for a monthly rate of $\dfrac{3.4}{1200}$ (or $0.002833\ldots$)
and an exponent of $144$. Using $3.4\%$ compounded annually (giving
\$$2987.28$) earns no marks.}

\medskip

\textbf{(c)} Bank A \hfill \textbf{A1}

\emph{FT: award A1 for the bank with the greater value from the
candidate's own answers to (a) and (b).}

\medskip

\textbf{(d)} $2000 \times 1.035^{\,n} \ge 4000$ \hfill \textbf{M1}

After $20$ years: \$$3979.58$ \ ($< 4000$); \quad after $21$ years:
\$$4118.86$ \ ($\ge 4000$)

$21$ years \hfill \textbf{A1}

\emph{Accept $1.035^{\,n} \ge 2$ solved on the GDC or with logarithms:
$n \ge 20.15\ldots$, so $21$ years. Accept the TVM solver for $N$ with
$PV = -2000$, $FV = 4000$.}

\emph{Accept a table search: full marks for showing both adjacent
years and stating $21$. A single row (after $21$ years the value is
\$$4118.86$, so $21$) still earns full marks, with the feedback note
that the adjacent year $20$ is what justifies ``first''. See
markscheme-spec §18.}

\emph{$20.1$ is not an acceptable final answer --- the answer must be
a whole number of years. Rounding down to $20$ earns the M1 only.}

\end{enumerate}
\end{document}
